\vspace{0.5em}
\textbf{Аннотация.} В работе предложен метод выявления вредоносной активности в информационных системах на основе атрибутивных метаграфов — математической модели, сочетающей структурные и семантические характеристики компонентов системы. Подход позволяет формализовать тактики атакующих, описанные в MITRE ATT\&CK, как паттерны на метаграфе и сопоставлять их с реальными событиями, зарегистрированными в SIEM-системах. Экспериментальная оценка на синтетических и реальных данных показала снижение числа ложных срабатываний на 19–23\,\% по сравнению с традиционными сигнатурными методами при сохранении полноты обнаружения выше 92\,\%. Метод поддерживает инкрементальное обновление модели при появлении новых техник в ATT\&CK.

\vspace{0.5em}
\textbf{Ключевые слова:} атрибутивные метаграфы, обнаружение кибератак, MITRE ATT\&CK, SIEM, метаграфовое моделирование, информационная безопасность.
\vspace{1em}